\exercisegroup

\begin{enumerate}
\def\labelenumi{\arabic{enumi})}
\item
  Show that the completion of a Hausdorff barrelled space is barrelled.
\item
  Let E be a vector space over R or C. Show that E, with the finest locally convex topology on E (II, p.~25) is barrelled. Deduce from this examples of barrelled spaces which are not metrizable and are not Baire spaces.
\item
  Let \(E\) be a Hausdorff locally convex space with a countably infinite base \((a_{n})\) .
\end{enumerate}

\begin{enumerate}
\def\labelenumi{\alph{enumi})}
\item
  Show that E admits a countable, topologically independent base \((e_n)\) (using the fact that every line in E has a topological complement, define the \(e_n\) by induction).
\item
  Show that, for E to be barrelled, it is necessary and sufficient that the topology T of E is identical with the finest locally convex topology on E (observe that the convex balanced envelope of every sequence \((\lambda_{n}e_{n})\) is closed in E). In particular, if T is metrizable, E is not barrelled (cf.~exerc. 2).
\end{enumerate}

\begin{enumerate}
\def\labelenumi{\arabic{enumi})}
\setcounter{enumi}{3}
\item
  Let E be a Banach space in which there exists an infinite algebraically independent sequence \((a_{n})\) which is total in E (for example the space \(\ell^{1}(\mathbf{N})\) (I, p.~4)). Let B be a base of E containing the \(a_{n}\) ; we know (II, p.~80, exerc. 24) that B is not countable. Let \((e_{n})\) be a sequence of distinct elements of B, and distinct from the \(a_{n}\) , and let C be the complement of the set of the \(e_{n}\) in B. Let \(F_{n}\) be the vector subspace of E generated by C and the \(e_{k}\) for indices \(k \leqslant n\) ; E is the union of the \(F_{n}\) . Let S be the unit ball in E; show that there exists an index n such that \(S \cap F_{n}\) , is not a first category set. Deduce that for this value of n, \(F_{n}\) is a metrizable, non-complete Baire space.
\item
  Give an example of a locally convex space which is a complete, Hausdorff Baire space, but is not metrizable (cf.~GT, IX, § 5, exerc. 16).
\item
  A locally convex space \(\mathbf{E}\) is said to be relatively bounded if there exists a bounded barrel in \(\mathbf{E}\) .
\end{enumerate}

\begin{enumerate}
\def\labelenumi{\alph{enumi})}
\item
  In order that E be relatively bounded, it is necessary and sufficient that the topology of E is coarser than a topology defined by a semi-norm. Then there exists a base for the canonical bornology of E consisting of barrels.
\item
  For E to be bornological and relatively bounded, it is necessary and sufficient that the topology of E is the lower bound of a family of normed space topologies on E (cf.~III, p.~40, exerc. 1). Further, in order that there exist also a countable base for the canonical bornology of E, it is necessary and sufficient that the topology of E is the lower bound of a sequence of normed space topologies.
\end{enumerate}

\begin{enumerate}
\def\labelenumi{\arabic{enumi})}
\setcounter{enumi}{6}
\item
  A locally convex space E is said to be infra-barrelled if every barrel of E which is bornivorous (III, p.~39, exerc. 11) is a neighbourhood of 0 in E. Every bornological space is infra-barrelled; every barrelled space is infrabarrelled. Show that the completion of a Hausdorff infrabarrelled space is barrelled (use the fact that in a Hausdorff locally convex space E, each barrel absorbs every convex, balanced, bounded and semi-complete subset of E).
\item
  Let \((\mathrm{E}_{\iota})_{\iota\in I}\) be a family of infrabarrelled spaces, and for every \(\iota\in I\) , let \(f_{\iota}\) be a linear mapping from \(\mathrm{E}_{\iota}\) into a vector space E. Show that the space E, with the finest locally convex topology for which the \(f_{\iota}\) are continuous, is infrabarrelled. In particular, every quotient space of an infrabarrelled space is infrabarrelled; every topological direct sum of infrabarrelled spaces is infrabarrelled.
\item
  Let I be an uncountably infinite set; on the direct sum vector space \(E = R^{(I)}\) , consider, on the one hand, the finest locally convex topology T, and on the other hand, the topology \(T_{0}\) defined in I, p.~24, exerc. 14, which is locally convex; we know that T and \(T_{0}\) are distinct (II, p.~75, exerc. 11) and that E with \(T_{0}\) is complete (GT, III, § 3, exerc. 10). Show that the bounded sets in E are the same for T and \(T_{0}\) (III, p.~39, exerc. 10) and that E with \(T_{0}\) is not barrelled (observe that the set T of all \(x = (\xi_{\mathfrak{u}}) \in \mathbf{E}\) such that \(\sum_{\mathfrak{u} \in I} |\xi_{\mathfrak{u}}| \leqslant 1\) is a barrel and use exerc. 11 of II, p.~75).
\item
  Show that an infrabarrelled space in which every closed convex balanced and bounded subset is semi-complete is a barrelled space.
\item
  Let \(\mathbf{E}\) be an infrabarrelled space, \(\mathbf{F}\) a locally convex space. Show that every subset of \(\mathcal{L}(\mathbf{E};\mathbf{F})\) which is bounded for the topology of bounded convergence is equicontinuous.
\end{enumerate}

¶ 12) a) Let E be a locally convex space, \((A_n)\) an increasing sequence of convex balanced sets in E such that \(A = \bigcup_{n} A_n\) is absorbent. Let \((W_n)\) be a decreasing sequence of convex balanced

neighbourhoods of 0; then the convex balanced envelope V of the \(\mathrm{W}_n\cap \mathrm{A}_n\) is absorbent; if E is barrelled, \(\overline{\mathbf{V}}\) is a neighbourhood of 0.

\begin{enumerate}
\def\labelenumi{\alph{enumi})}
\setcounter{enumi}{1}
\item
  Let \(\mathfrak{F}\) be a filter on E; suppose that for every \(n\) , there exists a set \(\mathrm{M}_n \in \mathfrak{F}\) such that \((\mathrm{M}_n + \mathrm{W}_n) \cap \mathrm{A}_{2n} = \emptyset\) . Let \(\mathrm{V}_n\) be the convex balanced envelope of the \(\mathrm{W}_k \cap \mathrm{A}_k\) for \(k \leqslant n - 1\) and of \(\mathrm{W}_n\) in such a way that \(\mathrm{V}_n\) is a neighbourhood of 0 and that we have \(\frac{1}{2}\overline{\mathrm{V}} \subset \mathrm{V}_n\) for all \(n\) . Show that \((\mathrm{M}_n + \mathrm{V}_n) \cap \mathrm{A}_n = \emptyset\) for all \(n\) .
\item
  Deduce from a) and b) that if E is barrelled and if \(\mathfrak{F}\) is a Cauchy filter on E, then there exists an integer N such that, for all \(\mathbf{M} \in \mathfrak{F}\) and every neighbourhood W of 0 in E, \(\mathbf{M} + \mathbf{W}\) meets \(\mathrm{A}_{\mathbf{N}}\) . (Argue by contradiction; with the notations of b), consider a set \(\mathbf{M} \in \mathfrak{F}\) with small order \(\frac{1}{2}\overline{\mathrm{V}}\) .)
\end{enumerate}

¶ 13) a) Let E be a barrelled space, \((C_n)\) an increasing sequence of convex, balanced sets such that \(E = \bigcup_{n} C_n\) . Let U be a convex, balanced and absorbent set such that for every n,

\(U \cap C_{n}\) is closed in \(C_{n}\) . Show that U is a neighbourhood of 0 in E. (Show that \(\overline{U} \subset 2U\) , by considering a filter \(\mathfrak{F}\) on U converging to a point \(x \in E\) and applying exerc. 12, c)).

\begin{enumerate}
\def\labelenumi{\alph{enumi})}
\setcounter{enumi}{1}
\tightlist
\item
  Let \(\mathrm{E}\) be a barrelled space, \((\mathrm{E}_n)\) an increasing sequence of subspaces of \(\mathbf{E}\) such that \(\mathrm{E} = \bigcup_{n} \mathrm{E}_n\) .
\end{enumerate}

Show that if \(\mathbf{U}\) is a subset of \(\mathbf{E}\) such that \(\mathbf{U} \cap \mathbf{E}_n\) is a barrel in \(\mathbf{E}_n\) for every \(n\) , then \(\mathbf{U}\) is a neighbourhood of 0 in \(\mathbf{E}\) . In particular, \(\mathbf{E}\) is the strict inductive limit (II, p.~33) of the sequence \((\mathbf{E}_n)\) .

¶ 14) a) Let E be a Hausdorff locally convex space, L a subspace of E with finite codimension, and T a barrel in L. Show that there exists a barrel T' in E such that T' ∩ L = T (show that we can take for T' the sum of the closure T̅ of T in E and of a finite dimensional compact convex set).

\begin{enumerate}
\def\labelenumi{\alph{enumi})}
\setcounter{enumi}{1}
\tightlist
\item
  Let E be a barrelled space, L a subspace of E, which has a complement having a countable basis. Show that L is barrelled (use \(a\) ) and exerc. 13, b)).
\end{enumerate}

* c) Let E be a Hausdorff locally convex space; its completion \(\hat{\mathbf{E}}\) can be identified with a closed subspace of a barrelled space F, which is the product of a family of Fréchet spaces (II, p.~5, prop. 3 and IV, p.~14, corollary). Let \((e_{\alpha})_{\alpha \in \mathbf{A}}\) be the basis of a complement in F of the subspace E, and let \(\mathrm{H}_{\alpha}\) be the hyperplane in F generated by E and the \(e_{\beta}\) for indices \(\beta \neq \alpha\) ; by \(b\) ), \(\mathrm{H}_{\alpha}\) is a barrelled space. For every \(x \in \mathrm{E}\) , let \(u(x)\) be the point of the barrelled space \(\mathrm{G} = \prod_{\alpha \in \mathbf{A}} \mathrm{H}_{\alpha}\) (sub-

space of \(F^{A}\) ) all whose coordinates are equal to x; u is an isomorphism from E onto the subspace \(\Delta \cap G\) , where \(\Delta\) is the diagonal in \(F^{A}\) . Show that \(u(E) = \Delta \cap G\) is closed in G, and consequently that every Hausdorff locally convex space is isomorphic to a closed subspace of a Hausdorff barrelled space. *

\begin{enumerate}
\def\labelenumi{\arabic{enumi})}
\setcounter{enumi}{14}
\tightlist
\item
  Let E be a Hausdorff barrelled (resp. infrabarrelled) space, and \(\hat{E}\) be its completion. Show that every subspace F of \(\hat{E}\) which contains E is barrelled (resp. infrabarrelled) (cf.~III, p.~24, cor. and IV, p.~52, exerc. 1).
\end{enumerate}

¶ * 16) Let \((E_{\iota})_{\iota\in I}\) be an uncountable family of Hausdorff barrelled spaces, none of which are the point 0, and let \(E = \prod_{\iota \in I} E_{\iota}\) ; then \(E\) is barrelled (IV, p.~14, corollary). Let \(G\) be the subspace of E consisting of all points \((x_{\iota})\) such that \(x_{\iota}=0\) except for a countable number of indices. Every sequence of points of G which converges in E has a limit belonging to G, but G is dense in E.

\begin{enumerate}
\def\labelenumi{\alph{enumi})}
\tightlist
\item
  Show that every subset \(\mathbf{M}\) of \(\mathrm{G}' = \mathrm{E}'\) , which is bounded for \(\sigma(\mathrm{E}', \mathrm{G})\) is contained in a finite product \(\prod_{\iota \in \mathrm{H}} \mathrm{E}_{\iota}'\) (IV, p.~12, prop. 13), where \(\mathrm{H}\) is a finite subset of I; consequently \(\mathbf{M}\) is
\end{enumerate}

bounded for \(\sigma(E', E)\) . Deduce from this that \(G\) is barrelled.

\begin{enumerate}
\def\labelenumi{\alph{enumi})}
\setcounter{enumi}{1}
\tightlist
\item
  Let F be a subspace of E such that \(G \subset F \subset E\) and such that G is a hyperplane (everywhere dense) in F; F is barrelled (exerc. 15). Show that F is not bornological. (Argue by reductio ad absurdum; if there were a convex, balanced and bounded set A in F such that G is an everywhere dense hyperplane in the normed space \(F_{A}\) (III, p.~7), then there would exist a sequence of points of G converging to a point of F not belonging to G) (cf.~IV, p.~52, exerc. 2). *
\end{enumerate}

\begin{enumerate}
\def\labelenumi{\arabic{enumi})}
\setcounter{enumi}{16}
\tightlist
\item
  \begin{enumerate}
  \def\labelenumii{\alph{enumii})}
  \tightlist
  \item
    Let E be a Hausdorff locally convex space, L a vector subspace of E of finite codimension, and T a bornivorous barrel in L. Show that there exists a bornivorous barrel \(T'\) in E such that \(T' \cap L = T\) . (Reduce to the case where L is a hyperplane in E. Let \(E_{0}\) be the bornological space associated with E (III, p.~40, exerc. 1), \(L_{0}\) the hyperplane L with the topology induced by that of \(E_{0}\) ; observe that T is a neighbourhood of 0 in \(L_{0}\) and consider the following two cases: that \(L_{0}\) is dense in \(E_{0}\) , or is closed in \(E_{0}\) ; show that for \(T'\) we can take the closure \(\overline{T}\) of T in E or the sum of T and a compact convex set of dimension 1).
  \end{enumerate}
\end{enumerate}

\begin{enumerate}
\def\labelenumi{\alph{enumi})}
\setcounter{enumi}{1}
\tightlist
\item
  Let E be an infrabarrelled space, L a vector subspace of E of finite codimension. Deduce from a) that L is infrabarrelled (cf.~IV, p.~64, exerc. 11).
\end{enumerate}

¶ 18) Let E be a strict inductive limit space of an increasing sequence ( \(E_{n}\) ) of locally convex metrizable subspaces (II, p.~33), and let F be a vector subspace of E such that every point of E is a limit point of a sequence of points of F.

\begin{enumerate}
\def\labelenumi{\alph{enumi})}
\item
  If \(\overline{\mathbf{E}}_n\) is the closure of \(\mathbf{E}_n\) in \(\mathbf{E}\) , then \(\mathbf{E}\) is the strict inductive limit of the sequence \((\overline{\mathbf{E}}_n)\) . Let \(\mathbf{F}_n\) be the closure of \(\mathbf{F} \cap \overline{\mathbf{E}}_n\) in \(\mathbf{E}\) . Show that \(\mathbf{E}\) is the union of the increasing sequence of subspaces \(\mathbf{F}_n\) .
\item
  Suppose E is barrelled. Show that F is bornological. (Let u be a linear mapping from F into a Banach space G which transforms every bounded subset of F into a bounded subset of G. Show that there exists a linear mapping from E into G, whose restriction to F is equal to u, and whose restriction to each \(F_{n}\) is continuous. Finally use exerc. 13, b) of III, p.~44.)
\end{enumerate}

\begin{enumerate}
\def\labelenumi{\arabic{enumi})}
\setcounter{enumi}{18}
\tightlist
\item
  A Hausdorff locally convex space E is said to be ultrabornological if every convex subset of E which absorbs all the convex, balanced, bounded and semi-complete subsets of E, is a neighbourhood of 0 in E.
\end{enumerate}

\begin{enumerate}
\def\labelenumi{\alph{enumi})}
\item
  Show that every ultrabornological space is both bornological and barrelled.
\item
  Let E be a Hausdorff locally convex space such that the closed, convex, balanced envelope of the set of points of every sequence tending to 0 is semi-complete. Show that if E is bornological then it is ultrabornological. In particular every bornological and quasi-complete space is ultrabornological; every Fréchet space is ultrabornological.
\item
  Let \((\mathrm{E}_{\alpha})\) be a directed increasing family of vector subspaces of a vector space E such that, E is the union of the \(\mathrm{E}_{\alpha}\) . Let \(\mathcal{T}_{\alpha}\) be a locally convex topology on \(\mathrm{E}_{\alpha}\) for every \(\alpha\) , and let \(\mathcal{T}\) be the finest locally convex topology for which the canonical injections from \(\mathrm{E}_{\alpha}\) into E are continuous. Suppose that \(\mathcal{T}\) is Hausdorff and that, for every \(\alpha\) , the topology on \(\mathrm{E}_{\alpha}\) induced by \(\mathcal{T}\) is \(\mathcal{T}_{\alpha}\) . Show that if each of the spaces \(\mathrm{E}_{\alpha}\) is ultrabornological, then E with \(\mathcal{T}\) is ultrabornological. d) Show that every finite product of ultrabornological spaces is ultrabornological; deduce that every topological direct sum of ultrabornological spaces is ultrabornological.
\item
  Show that the product space \(\mathrm{E} = \prod_{n=0}^{\infty} \mathrm{E}_n\) of an infinite sequence of ultrabornological
\end{enumerate}

spaces is ultrabornological. (Let A be a convex subset of E which absorbs every convex, balanced, bounded and semi-complete subset of E. Show that if A were not a neighbourhood of 0 in E, then there would have existed a sequence \((x_{n})\) in \(\mathbb{C}\) A such that \(x_{n}\) has its first \(n - 1\) coordinates zero, but is \(\neq 0\) . Next observe that the closed convex balanced envelope of the set of points of such a sequence is identical to the set of points \(\sum_{n = 0}^{\infty}\lambda_n x_n\) , where \(\sum_{n = 0}^{\infty}|\lambda_n| \leqslant 1\) , and

that this envelope is a semi-complete set.)

\begin{enumerate}
\def\labelenumi{\arabic{enumi})}
\setcounter{enumi}{19}
\item
  Show that, for a Hausdorff locally convex space \(E\) to be ultrabornological it is necessary and sufficient that it is the inductive limit of a family of Banach spaces. (To see that the condition is necessary, consider the convex, balanced, bounded and semi-complete sets \(B\) in \(E\) , and the spaces \(E_B\) . To see that it is sufficient, observe that if \(E\) is the inductive limit of a family of Banach spaces \(E_\alpha\) , we can assume that the \(E_\alpha\) are (algebraically) subspaces of \(E\) ; if \(V\) is a convex set in \(E\) which absorbs the convex, balanced, bounded and semi-complete subsets of \(E\) , show that \(V\) absorbs each ball \(B_\alpha\) of \(E_\alpha\) (argue by reductio ad absurdum); if \(V\) does not absorb \(B_\alpha\) , then it does not absorb a sequence \((x_n)\) of points of \(B_\alpha\) , tending to 0 in \(E\) ; then use the fact that in a Banach space, the closed, convex envelope of a compact set is compact.)
\item
  Show that if E is a Hausdorff locally convex semi-complete space, then the bornological space associated with E (III, p.~40, exerc. 1) is ultrabornological.
\end{enumerate}

¶ 22) Let E be an infinite dimensional Banach space satisfying the first axiom of countability.

\begin{enumerate}
\def\labelenumi{\alph{enumi})}
\item
  Show that the set \(\mathcal{K}\) of all compact, convex and balanced subsets A of E such that \(\mathrm{E_A}\) is infinite dimensional, is infinite and has a cardinality \(\leqslant 2^{\mathrm{Card}(\mathbf{N})}\) (GT, IX, § 5, exerc. 17). For every \(x_0 \in \mathrm{E}\) and every \(\mathrm{A} \in \mathcal{K}\) , the set \(x_0 + \mathrm{A}\) contains a free subset of cardinality \(2^{\mathrm{Card}(\mathbf{N})}\) (II, p.~80, exerc. 24, c)).
\item
  Let \(x_0 \neq 0\) be in E. Show that there exists a family \((y_{\mathrm{A}})_{\mathrm{A} \in \mathcal{K}}\) such that \(x_0\) and the \(y_{\mathrm{A}}\) form a free family and that we have \(y_{\mathrm{A}} \in x_0 + \mathrm{A}\) for all \(\mathrm{A} \in \mathcal{K}\) (well order \(\mathcal{K}\) and argue by transfinite induction, using \(a\) ).
\item
  Let \(f \in \mathrm{E}^{*}\) be a linear form such that \(f(x_0) = 1\) and \(f(y_{\mathrm{A}}) = 0\) for all \(\mathrm{A} \in \mathcal{K}\) , and let \(\mathrm{H} = f^{-1}(0)\) . Show that a subset \(\mathbf{M}\) of \(\mathrm{H}\) which is convex, balanced and semi-complete is necessarily finite dimensional (observe that if not, \(\mathbf{M}\) will contain an infinite dimensional compact convex and balanced set \(\mathrm{A}\) ; hence \(y_{\mathrm{A}}\) will belong to \(\mathrm{H} \cap (x_0 + \mathbf{M})\) ).
\item
  Show that H with the topology induced by that of E is not ultrabornological, in spite of being bornological and barrelled (III, p.~44, exerc. 14). (By using c), show that if H were ultrabornological, its topology would have been the finest locally convex topology, and deduce a contradiction.)
\end{enumerate}

